% !TEX program = xelatex
% Example: Two-column frame layout in Beamer
% Left column holds text; right column holds a simple TikZ diagram.
\documentclass{beamer}
\usetheme{default}
\usepackage{tikz}

\title{Column Layouts}
\author{Dr.~Niels Bohr}
\institute{Institute of Physics}
\date{2025}

\begin{document}

\begin{frame}
  \titlepage
\end{frame}

\begin{frame}{Two-Column Layout}
  \begin{columns}[T,onlytextwidth]
    \begin{column}{0.48\textwidth}
      \textbf{Left column: text}
      \begin{itemize}
        \item Columns let you place content side by side.
        \item Useful for text-plus-diagram slides.
        \item Use \texttt{[T]} to top-align both columns.
      \end{itemize}
    \end{column}
    \begin{column}{0.48\textwidth}
      \textbf{Right column: diagram}
      \begin{center}
        \begin{tikzpicture}
          \node[draw, rounded corners, fill=blue!15,
                minimum width=2.5cm, minimum height=1cm] (a) {Input};
          \node[draw, rounded corners, fill=green!15,
                minimum width=2.5cm, minimum height=1cm,
                below=0.6cm of a] (b) {Process};
          \node[draw, rounded corners, fill=red!15,
                minimum width=2.5cm, minimum height=1cm,
                below=0.6cm of b] (c) {Output};
          \draw[-stealth] (a) -- (b);
          \draw[-stealth] (b) -- (c);
        \end{tikzpicture}
      \end{center}
    \end{column}
  \end{columns}
\end{frame}

\end{document}
